% ============================================================
%  INTRODUCTION GÉNÉRALE
% ============================================================
\chapter*{Introduction Générale}
\addcontentsline{toc}{chapter}{Introduction Générale}
\markboth{Introduction Générale}{Introduction Générale}

\section*{Contexte général}

L'essor de l'intelligence artificielle et de l'apprentissage automatique a profondément
transformé la manière dont les organisations gèrent leurs infrastructures informatiques.
Dans un contexte où la disponibilité des systèmes est devenue un enjeu stratégique majeur,
la maintenance prédictive s'impose comme une approche incontournable pour anticiper les
défaillances matérielles avant qu'elles ne surviennent.

Les parcs informatiques modernes génèrent en permanence des volumes considérables de données
de surveillance : utilisation du processeur, consommation mémoire, état des disques durs,
températures des composants. Ces données, correctement analysées, constituent une source
d'information précieuse pour prédire les pannes et planifier les interventions de maintenance
de manière proactive.

La transformation numérique des entreprises a conduit à une dépendance croissante vis-à-vis
des systèmes informatiques. Une panne non anticipée peut entraîner des pertes financières
considérables, une interruption des services critiques, et une dégradation de la confiance
des utilisateurs. Face à cette réalité, les technologies d'apprentissage automatique offrent
des perspectives prometteuses : les algorithmes de classification, les réseaux de neurones
récurrents et les modèles de détection d'anomalies permettent d'identifier les signes
précurseurs de défaillance bien avant que la panne ne se produise.

\section*{Problématique}

La maintenance réactive, qui consiste à intervenir uniquement après la survenue d'une panne,
engendre des coûts importants : temps d'arrêt non planifiés, perte de productivité, urgences
techniques et remplacement précipité de matériel. La question centrale de ce projet est :

\begin{infobox}[Problématique centrale]
Comment concevoir et implémenter un système intelligent capable de surveiller en temps réel
un parc informatique, d'analyser les tendances des métriques matérielles, et de prédire les
pannes avant qu'elles ne se produisent, afin de permettre une maintenance proactive et
réduire les temps d'arrêt non planifiés ?
\end{infobox}

Cette problématique soulève plusieurs questions secondaires :
\begin{itemize}[leftmargin=1.5cm]
  \item Comment collecter de manière automatique et fiable les métriques système (CPU, RAM,
        disque, SMART) sur un parc hétérogène de machines Windows et Linux ?
  \item Quels algorithmes d'apprentissage automatique sont les plus adaptés à la prédiction
        de défaillances à partir de séries temporelles de métriques système ?
  \item Comment évaluer objectivement la qualité d'un assistant conversationnel spécialisé
        dans le domaine de la maintenance informatique ?
  \item Comment garantir la confidentialité des données d'infrastructure tout en offrant
        des capacités conversationnelles avancées ?
\end{itemize}

\section*{Objectifs du projet}

Ce projet vise à développer une plateforme complète de maintenance prédictive reposant sur
l'intelligence artificielle. Les objectifs principaux sont :

\begin{itemize}[leftmargin=1.5cm]
  \item Mettre en place un agent Python de collecte automatique des métriques système
        (via \texttt{psutil}) et des données SMART (via \texttt{smartctl}) sur chaque
        machine surveillée ;
  \item Concevoir un pipeline d'extraction de 63 caractéristiques statistiques à partir
        des données historiques stockées dans PostgreSQL ;
  \item Entraîner et déployer deux modèles complémentaires : un Random Forest pour les
        métriques système et un LSTM entraîné sur le dataset Backblaze pour les données
        SMART ;
  \item Développer un système d'alertes automatiques par email (Nodemailer/SMTP) pour
        les niveaux de risque HIGH et CRITICAL ;
  \item Créer un tableau de bord React interactif pour la visualisation des métriques,
        des prédictions à 7/14/30 jours, et la gestion des alertes ;
  \item Intégrer un assistant conversationnel hybride basé sur LLaMA 3.2:1b (Ollama)
        fonctionnant entièrement en local, sans transmission de données vers le cloud ;
  \item Évaluer les performances du chatbot avec les métriques ROUGE-1 et ROUGE-2 sur
        un dataset de 20 questions-réponses de référence.
\end{itemize}

\section*{Organisation du rapport}

Ce rapport est organisé en quatre chapitres principaux :

\begin{description}[leftmargin=2cm, style=nextline]
  \item[\textbf{Chapitre 1 — Vue d'ensemble du projet :}] présente le contexte, la
  problématique, les objectifs, la solution proposée et la méthodologie Agile adoptée.

  \item[\textbf{Chapitre 2 — État de l'art :}] passe en revue les fondements théoriques
  de l'apprentissage automatique appliqué à la maintenance prédictive, les modèles LSTM,
  les LLM, et les métriques d'évaluation des systèmes conversationnels.

  \item[\textbf{Chapitre 3 — Conception du système :}] détaille l'architecture
  microservices, la modélisation UML, les besoins fonctionnels et non fonctionnels,
  et les flux de données entre les composants.

  \item[\textbf{Chapitre 4 — Implémentation technique :}] décrit les technologies
  utilisées, l'implémentation de chaque module, les résultats d'évaluation quantitatifs,
  et les captures d'écran de l'interface.
\end{description}

Le rapport se termine par une conclusion générale récapitulant les réalisations et les
perspectives d'évolution du système.
